\begin{lemma}[Terminal PE coloring and closure]
Let \(\mathcal E:\noderef{ExtendedTashkinovTree}(G,k,e_0,\rho,T,\varphi)\)
be an ETT in a finite loopless multigraph. Put \(n=\mathcal E.\mathit{rounds}\),
\(T_n=\mathcal E.\mathit{series.data.tree}(\operatorname{last}(n).\mathit{castSucc})\),
and \(F=\mathcal E.\mathit{series.data.tree}(\operatorname{last}(n+1))\).
If its terminal extension kind is PE, then \(\varphi\) is proper, leaves
exactly \(e_0\) uncolored, and \(F\) is a
\noderef{TashkinovClosure} of \(T_n\) under \(\varphi\).
\end{lemma}
\begin{proof}
The positivity field of \noderef{ExtendedTashkinovTree} gives \(n\geq1\).
Take \(q=n-1\) as an index in \(\operatorname{Fin}(n)\); its successor
is the terminal stored index \(n\), and its double successor is the final
tree index \(n+1\). The certificate supplied by \noderef{TashkinovSeries}
at index \(n\) has one of the four forms in \noderef{TashkinovSeriesStep}.
It cannot be initial, whose index is zero. The SE and RE forms record those
respective kinds, contradicting the assumed PE kind. In the remaining PE
form, the current-properness and current-one-uncolored fields say that the
coloring at index \(n\) is proper and leaves exactly \(e_0\) uncolored.
Its closed-extension field states that the tree at index \(n+1\) is a
closure of the tree at index \(n\), under that current coloring.
The generating-coloring equality of \noderef{ExtendedTashkinovTree}
identifies the current coloring with \(\varphi\). Substituting this equality
and the two index identities gives all three conclusions.
\end{proof}
